\section{A Research Agenda: What Would Actually Move the Needle}
\label{sec:future}

If the distance estimate of Section~\ref{sec:distance} is right, then the
highest-leverage work is not bigger models but better \emph{infrastructure
and protocols}. We propose a concrete agenda, ordered by leverage.

\subsection{A public, grounded canine corpus}
The field lacks what BEANS and BirdSet provide for general bioacoustics: a
standardized, public, multimodal canine corpus with agreed
labels~\cite{hagiwara2023beans}. The ideal corpus pairs synchronized audio,
video (DogFACS-codeable face and posture), context metadata, and---crucially%
---\emph{consequent behavior}, so that intent hypotheses can be checked. The
soundboard citizen-science network, with thousands of instrumented homes,
is the natural collection substrate~\cite{bastos2024plos}. Building this
corpus with disclosed protocols and inter-annotator agreement would do more
for the field than any single model release.

\subsection{Calibrated, multimodal affect models}
On the modeling side, the priority is fusion and calibration: combine bark,
growl, face, and posture into a single state estimate with honest
uncertainty, trained and evaluated on the public corpus. Self-supervised
encoders (AVES, NatureLM-audio, Perch) supply the
backbone~\cite{hagiwara2023aves,robinson2024naturelm,ghani2023perch}; the
contribution is the canine-specific grounding and the calibration layer
that makes the output safe for welfare use.

\subsection{Closed-loop intent validation}
The methodological frontier is the validation loop: infer an intent, act on
it (or emit a candidate signal), and test the dog's response under blind
conditions. The soundboard and CHAT protocols are the
templates~\cite{bastos2024scirep,google2025dolphingemma}. A home-adapted,
owner-blind protocol that scores intent inference against the dog's
subsequent behavior would supply the error signal the field lacks
(Section~\ref{sec:eval}) and turn Rung 3 from hypothesis into measurement.

\subsection{Structure discovery, with humility}
Applying sequence-structure methods to bark corpora is worth doing---but
with the explicit prior that canine displays are graded and may yield no
clean combinatorial alphabet. The CETI and DolphinGemma toolkits are
available~\cite{sharma2024whale,google2025dolphingemma}; the discipline is to
report structure as form, not meaning, and to heed the acoustic-similarity
critique~\cite{yovel2025critique}. A negative result---no discrete
inventory---would itself be valuable, sharpening the semantic-gap claim.

\subsection{Engineered dialogue, not extracted dialogue}
The most promising path to anything resembling conversation is
\emph{engineered}: expand the soundboard paradigm into a richer, negotiated
symbol system, with the dog as partner and the behavior as check. This is
Rung-3 work that may eventually look like dialogue, without pretending that
spontaneous barks contain sentences. It is also the ethically preferable
path (Section~\ref{sec:ethics}).

\subsection{What would falsify our thesis}
A good survey states what would change its mind. Our central claim---that
Rungs 4--5 are ill-posed for natural canine communication---would be
undermined by: (i) a reproducible demonstration of a productive, open
referential vocabulary in spontaneous dog signals, validated by playback;
or (ii) evidence of recursive, generative combination in natural canine
displays (not trained soundboard sequences). Absent such evidence, the
ladder stands. We invite the community to try to falsify it; that is how
the distance estimate should shrink.